%% LyX 2.4.4 created this file.  For more info, see https://www.lyx.org/.
%% Do not edit unless you really know what you are doing.
\documentclass[english]{article}
\usepackage[T1]{fontenc}
\usepackage[latin9]{inputenc}
\usepackage{units}
\usepackage{enumitem}
\usepackage{amsmath}
\usepackage{amsthm}
\usepackage{amssymb}
\usepackage{geometry}
\geometry{verbose,tmargin=1in,bmargin=1in,lmargin=1in,rmargin=1in}

\makeatletter
%%%%%%%%%%%%%%%%%%%%%%%%%%%%%% Textclass specific LaTeX commands.
\numberwithin{equation}{section}
\numberwithin{figure}{section}
\newlength{\lyxlabelwidth}      % auxiliary length 

%%%%%%%%%%%%%%%%%%%%%%%%%%%%%% User specified LaTeX commands.
\usepackage{fancyhdr}
\usepackage{lastpage}
\pagestyle{fancy}
\usepackage{enumitem}
\usepackage{ifthen}
\everymath=\expandafter{\the\everymath\displaystyle}
%\DeclareMathSizes{10}{10}{10}{10}
\usepackage{multicol}

\usepackage{tikz}
\usepackage{tikz-3dplot}
\usetikzlibrary{calc,arrows}

\usetikzlibrary{patterns}
\usetikzlibrary{plotmarks}
\usepackage{pgfplots}
\pgfplotsset{compat=newest}

\usepackage{calc}
\usepackage[nomessages]{fp}

\usepackage{datetime}
% Added by lyx2lyx
\setlength{\parskip}{\smallskipamount}
\setlength{\parindent}{0pt}

\makeatother

\usepackage{babel}
\begin{document}
\renewcommand{\author}{Dr.\,James Ashe}

\newcommand{\classNum}{336}

\newcommand{\classSec}{A}

\newcommand{\examType}{Final Exam}

\newcommand{\nameType}{Name:}

\renewcommand{\date}{\formatdate{4}{12}{2013}}

%%First page's header%%%%%%%%%%%%%%%%%%%%%%
\fancypagestyle{firststyle} %%header settings for first pagr
{
	\fancyhead[L]{MTH \classNum{}(\classSec{}) \author{}\\
	\textbf{\examType{}} ~\date{}}
	\fancyhead[C]{\textbf{\nameType}}
	\fancyhead[R]{}
	\setlength{\headsep}{14pt}
}
%%%%%%%%%%%%%%%%%%%%%%%%%%%%%%%%%%%%%%%%%%

%%Next page's header%%%%%%%%%%%%%%%%%%%%%%%
\setlist{nolistsep}
\lhead{MTH \classNum{}(\classSec{})} 
\chead{\textbf{\nameType}} 
\cfoot{\thepage{}~of~\pageref{LastPage}} 
\setlength{\headsep}{5pt} 
%%%%%%%%%%%%%%%%%%%%%%%%%%%%%%%%%%%%%%%%%%%

%%Various settings; see the preamble for other settings%%%%%
%\setenumerate[0]{leftmargin=12pt,itemindent=0pt,labelwidth=0pt}
\setlist{topsep=.5pt}
\renewcommand{\labelenumi}{\textbf{\arabic{enumi}.}}
\renewcommand{\labelenumii}{\textbf{\alph{enumii}.}}
\thispagestyle{firststyle}
%%%%%%%%%%%%%%%%%%%%%%%%%%%%%%%%%%%%%%%%%%

\newcommand{\xmax}{0}
\newcommand{\xmin}{-\xmax}
\newcommand{\xstep}{0}
\newcommand{\ymax}{0}
\newcommand{\ymin}{-\ymax}
\newcommand{\ystep}{0} 
\newcommand{\dspace}{\vspace{1cm}}
\newcommand{\xa}{0}
\newcommand{\xb}{0}
\newcommand{\ya}{1}
\newcommand{\yb}{1}

\newcommand{\xspace}{\vspace{.5cm}}
\newcommand{\Xspace}{\vspace{1cm}}

\newcommand{\vectors}[2]{\textbf{#1}_{1},\textbf{#1}_{2},\dots,\textbf{#1}_{#2}}
\newcommand{\vectorsN}[1]{\textbf{#1}_{1},\textbf{#1}_{2},\dots,\textbf{#1}_{n}}
\newcommand{\R}[1]{\mathbb{R}^{#1}}
\newcommand{\mc}[1]{[\textbf{(#1.)}]}

%\newcounter{probs}

%%for multiple choice
%\renewcommand{\labelenumi}{\textbf{\arabic{enumi}.}}
\renewcommand{\labelenumii}{\textbf{(\Alph{enumii}})}

\textbf{Work must be shown to receive credit. All problems are of
equal value.}
\begin{enumerate}[resume]
\item Let ${\vectors{v}{n+1}}$ be a set of $n+1$ vectors in $\mathbb{R^{\mbox{n}}}$.
The which of the following is true? Explain.

\begin{enumerate}[resume]
\item ${\vectors{v}{n+1}}$ is a basis for $\mathbb{R}^{\mbox{n}}$.
\item ${\vectors{v}{n+1}}$ is linearly independent.
\item ${\vectors{v}{n+1}}$ is linearly dependent .
\item The Span$(\vectors{v}{n+1})$ is not a subspace of $\mathbb{R}^{\mbox{n}}$.\Xspace
\end{enumerate}
\item If the reduced row echelon form of a system is $\begin{bmatrix}\begin{array}{rrr|r}
1 & 0 & 3 & 29\\
0 & 1 & 1 & 8\\
0 & 0 & 0 & 0
\end{array}\end{bmatrix}$. Then the system has,

\begin{enumerate}
\item a unique solution.
\end{enumerate}
\begin{enumerate}[resume]
\item no solution.
\item infinitely many solutions.
\item two solutions.\xspace
\end{enumerate}
\item If the \emph{reduced row echelon form} of the augmented matrix of
a system of linear equations has a leading coefficient 1 in the last
column then the system has:

\begin{enumerate}
\item a unique solution.
\end{enumerate}
\begin{enumerate}[resume]
\item no solution.
\item infinitely many solutions.
\item none of the above.\xspace
\end{enumerate}
\item If $\det A=6$ and $AB=\begin{bmatrix}\begin{array}{rrr}
-2 & 0 & 0\\
3 & -1 & 0\\
1 & 5 & 1
\end{array}\end{bmatrix}$ find $\det B$.\vfill
\item Let $A=\begin{bmatrix}\begin{array}{rrr}
1 & 2 & -1\\
-1 & 3 & 2\\
0 & 2 & 1
\end{array}\end{bmatrix}.$ Find the characteristic polynomial of $A$.\vfill
\end{enumerate}
\newpage
\begin{enumerate}[resume]
\item Let $A'$ be the reduced row echelon form of a matrix $A$. Then which
of the following is true.

\begin{enumerate}
\item $\det A=\det A'$
\end{enumerate}
\begin{enumerate}[resume]
\item $\det A=\pm\det A'$
\item $\det A=\lambda\det A'$ for some nonzero scalar $\lambda$.
\item $\det A$ and $\det A'$ are unrelated.\xspace
\end{enumerate}
\item Let $A$, $B$, and $C$ be three matrices such that $AB=CA$ and
$A$ is invertible. Then,

\begin{enumerate}
\item $\det B=\det C$
\end{enumerate}
\begin{enumerate}[resume]
\item $\det B=\pm\det C$
\item $\det B=-\det C$
\item $\det B$ and $\det C$ are unrelated. \xspace
\end{enumerate}
\item Let $A$ be an $m\times n$ real matrix. Then the function $T[\mathbf{x}]=A\mathbf{x}$
is 

\begin{enumerate}
\item a linear map from $\R{m}$ to $\R{n}$.
\end{enumerate}
\begin{enumerate}[resume]
\item a non-linear map from $\R{m}$ to $\R{n}$.
\item a linear map from $\R{n}$ to $\R{m}$.
\item not a function.\xspace
\end{enumerate}
\item Let $A=\begin{bmatrix}\begin{array}{rrr}
1 & 4 & 7\\
2 & 5 & 8\\
3 & 6 & 9
\end{array}\end{bmatrix}.$ Then which of the following is true.

\begin{enumerate}
\item $A$ is singular.
\end{enumerate}
\begin{enumerate}[resume]
\item $A$ is invertible and $A^{-1}=\begin{bmatrix}\begin{array}{rrr}
1 & \nicefrac{1}{4} & \nicefrac{1}{7}\\
\nicefrac{1}{2} & \nicefrac{1}{5} & \nicefrac{1}{8}\\
\nicefrac{1}{3} & \nicefrac{1}{6} & \nicefrac{1}{9}
\end{array}\end{bmatrix}$.
\item $\det A=1$.
\item $\det A=9$.\vfill
\end{enumerate}
\item Let $\mathbf{v}=(3,-2,4)$, $\mathbf{w}=(-3,2,-4)$, and $\mathbf{u}=(-6,4-8)$
be three vectors in $\R{n}$. Find the dimension of $\mbox{Span}(\mathbf{v},\mathbf{w},\mathbf{u})$
and explain your answer.\vfill
\end{enumerate}
\newpage
\begin{enumerate}[resume]
\item Let $A$ be an $n\times n$ real matrix. If the equation $A\mathbf{x}=\mathbf{0}$
has a trivial solution then, 

\begin{enumerate}
\item $A$ is invertible.
\end{enumerate}
\begin{enumerate}[resume]
\item $A$ is singular.
\item the system has infinitely many solutions.
\item none of the above.\xspace
\end{enumerate}
\item Let $A$ be a matrix with $\det A=3$. 

\begin{enumerate}
\item $\det(A^{-1})=\nicefrac{1}{3}$
\end{enumerate}
\begin{enumerate}[resume]
\item $\det(A^{-1})=\nicefrac{1}{3^{2}}$
\item $\det(A^{-1})$ is undefined because $A^{-1}$ might not exist.
\item $\det(A^{-1})$ and $\det A$ are unrelated.\xspace
\end{enumerate}
\item Let $A=\begin{bmatrix}\begin{array}{rrrr}
1 & 0 & 0 & 0\\
10 & 2 & 0 & 0\\
-7 & -1 & 3 & 0\\
25 & 5 & 11 & 5
\end{array}\end{bmatrix}$ find $\det A$.\vfill
\item Let $A=\begin{bmatrix}\begin{array}{rr}
5 & 2\\
7 & 3
\end{array}\end{bmatrix}$. Find $A^{-1}$ if it exists.\vfill
\item Find the inverse of the matrix $A=\begin{bmatrix}\begin{array}{rrrr}
4 & -3 & 0 & 0\\
-5 & 4 & 0 & 0\\
0 & 0 & 3 & -2\\
0 & 0 & 5 & -3
\end{array}\end{bmatrix}$ if it exists. \vfill
\end{enumerate}

\end{document}
